
\chapter{Residual Complexes}
\setcounter{section}{-1}
\begingroup
\setlength{\emergencystretch}{2em}

\section{Introduction}\label{sec:6.0}

In this chapter we return to the problem of constructing a functor $f^{!}$ for a morphism of finite type, which should reduce to $f^{\flat}$ for a finite morphism, and $f^{\sharp}$ for a smooth morphism. We ran into difficulty earlier \bookref{sec:3.8}{III,~\S8} because the derived category is not a local object — one cannot glue together elements of the derived category given locally. Now we overcome that difficulty in a special case by using residual complexes. The residual complex is a very special complex of quasi-coherent injective sheaves (see definition below) which is almost unique (it was called ``residue complex'' in the Edinburgh Congress talk \cite{9}).

Modulo some technicalities arising from possibly infinite Krull dimension, this is how they work: To each dualizing complex $R^\bullet\in\DD^+_{\mathrm{c}}(X)$ is associated functorially a residual complex $K^\bullet = E(R^\bullet)$, and $Q(K^\bullet)$ (which is the image of the complex $K^\bullet$ in $\DD^+_{\mathrm{c}}(X)$) is isomorphic to $R^\bullet$ in $\DD^+_{\mathrm{c}}(X)$. Thus we will define $f^{!}$ locally for a residual complex $K^\bullet$ by
\[
f^{!}(K^\bullet) = E(f^{!}(Q(K^\bullet))),
\]
where $f^{!}$ is the one we know for embeddable morphisms. Then since $f^{!}(K^\bullet)$ is an actual complex defined locally, we can glue to get a global $f^{!}(K^\bullet)$. We will give the statement and proof of the existence of $f^{!}$ in some detail, since this is an important step towards the general duality theorem. Later, after proving the duality theorem, we will pull ourselves up by our bootstraps to obtain a definition of $f^{!}$ for objects in the derived category $\DD_{\mathrm{c}}^{+}(X)$.

Once having $f^{!}$, we will define a trace map for residual complexes, which will be a map of graded sheaves
\[
\operatorname{Tr}_{f}\colon f_{*}f^{!}K^{\bullet}\longrightarrow K^{\bullet}.
\]
We will prove in the next chapter that if $f$ is a proper morphism, then $\Tr_f$ is a map of complexes (Residue Theorem).

\section{Residual Complexes}\label{sec:6.1}

Throughout this section, $X$ will denote a locally noetherian prescheme. If $x$ is a point of $X$, and if $I$ is an injective hull of $k(x)$ over the local ring $\sO_{X,x}$, we will denote by $J(x)$ the quasi-coherent injective $\sO_X$-module, which is the constant sheaf $I$ on $\overline{\{x\}}$, and $0$ elsewhere (notation of \bookref{sec:2.7}{II,~\S7}).

\begin{definition*}
A \emph{residual complex} on $X$ is a complex $K^\bullet$ of quasi-coherent injective $\sO_X$-modules, bounded below, with coherent cohomology sheaves, and such that there is an isomorphism
\[
\bigoplus_{p\in\mathbb{Z}} K^{p} \cong \bigoplus_{x\in X} J(x).
\]
\end{definition*}

\begin{example*}
If $X$ is a regular prescheme, then the Cousin complex of the structure sheaf $\sO_X$ is a residual complex for $X$ (see example at end of \bookref{sec:4.2}{IV,~\S2}).
\end{example*}

\begin{bookproposition}{1.1}\label{prop:6.1.1}
\begin{enumerate}
\item[a)] If $K^\bullet$ is a residual complex on $X$, then its image $Q(K^\bullet)\in\DD^+_{\mathrm{c}}(X)$ is a pointwise dualizing complex.
\item[b)] If $R^\bullet\in\DD^+_{\mathrm{c}}(X)$ is a pointwise dualizing complex on $X$, then $E(R^\bullet)$ is a residual complex on $X$. (Here $E$ is the notation of \bookref{sec:4.3}{IV,~\S3}, with respect to the filtration $Z^\bullet$ associated with $R^\bullet$ (cf.\ \bookref{sec:5.7}{V,~\S7} and \bookref{sec:5.8}{V,~\S8, Remark~4}).)
\item[c)] If $X$ admits a residual complex (or pointwise dualizing complex) with bounded cohomology, then there is a functorial isomorphism
\[
\varphi\colon R^{\bullet}\xrightarrow{\sim} QE(R^{\bullet})
\]
for pointwise dualizing complexes. Hence the functor
\[
E\colon \operatorname{Ptwdual}(X)\longrightarrow \operatorname{Res}(X)
\]
is an equivalence of the category of pointwise dualizing complexes of $\DD_{\mathrm{c}}^{+}(X)$ and the category of residual complexes (and morphisms of complexes). Its inverse is $Q$.\footnote{The category and boundedness qualifications are discussed in the editorial remark following the original remarks below.}
\end{enumerate}
\end{bookproposition}

\begin{proof}
a) The question is local, so we may assume $X=\Spec A$, where $A$ is a noetherian local ring. By \bookref{prop:5.3.4}{V,~3.4} we have only to check that there is an integer $d$ such that
\[
\operatorname{Ext}^{i}(k,K^{\bullet}) = \begin{cases} 0 & \text{for } i\neq d \\ k & \text{for } i = d \end{cases}
\]
where $k$ is the residue field of $A$. Since $k=k(x_0)$, where $x_0$ is the closed point of $X$, we have
\[
\Hom(k, J(x_0)) = k
\]
\[
\Hom(k, J(x)) = 0 \quad\text{for } x\neq x_0.
\]
The result now follows from the definition of a residual complex.

\noindent b) By the pointwise convergence of the spectral sequence \bookref{motif:4.1.G}{IV,~1.G} we see that $\mathcal{H}^{i}(R^{\bullet}) = \mathcal{H}^{i}(E(R^{\bullet}))$ for all $i$, and hence $E(R^{\bullet})$ is a complex which is bounded below and has coherent cohomology. By \bookref{motif:4.1.F}{IV,~1.F} the question of whether it is a residual complex is local, so we may assume $X$ is the spectrum of a local ring, in which case our result is \bookref{prop:5.7.3}{V,~7.3}.

\noindent c) This follows from \bookref{sec:5.8}{V,~\S8, Remark~5} and \bookref{prop:4.3.4}{IV,~3.4}.
\end{proof}

\par\medskip\noindent\textbf{Remarks.}
\begin{enumerate}
\item In particular, if $X$ admits a dualizing complex (and hence has finite Krull dimension), the functor
\[
E\colon \operatorname{Dual}(X)\longrightarrow \operatorname{Res}(X)
\]
is an equivalence of categories, with inverse $Q$.
\item I expect that the statement c) is false without the boundedness assumption. That is, there may be two non-isomorphic pointwise dualizing complexes $R^\bullet$ and $R^{\prime\bullet}$ such that $E(R^\bullet)=E(R^{\prime\bullet})$. In particular, there may not be a uniqueness theorem similar to \bookref{thm:5.3.1}{V,~3.1} for pointwise dualizing complexes.
\item It follows from the proposition that $X$ admits a residual complex if and only if it admits a pointwise dualizing complex, so the remarks of \bookref{sec:5.10}{V,~\S10} apply.
\end{enumerate}
\par\medskip

\subsection*{Editorial remark on Proposition 1.1: filtration and bounds}
For an equivalence with actual morphisms of complexes, fix a bounded-below codimension function $d$, put $Z^p=\{x\mid d(x)\geq p\}$, and use the categories $\operatorname{Ptwdual}_d(X)$ and $\operatorname{Res}_d(X)$ of objects with this function. Then $E$ and $Q$ give inverse equivalences. This is a qualification of the printed claim, not a change in what the author wrote.

For a DVR $A$ with fraction field $F$, let $C=(F\xrightarrow{\pi}F/A)$ in degrees $0,1$. The nonzero map $C\to C[1]$ with degree-zero component $\pi$ is null-homotopic: take $h^0=-\mathrm{id}_F$ and all other components zero. Thus $Q$ is not faithful when codimension functions vary, even for bounded residual complexes. In particular, a construction factoring through $Q$ cannot be naturally isomorphic to the actual identity functor on that broader category. This also affects unrestricted functoriality assertions in the following sections; an objectwise reconstruction alone does not settle them.

For bounded cohomology this follows from \bookref{sec:5.8}{V,~\S8, Remark~5} and \bookref{prop:4.3.4}{IV,~3.4}. For the stated bounded-below categories, use the Cousin equivalence for an admissible strictly exhaustive filtration. (Editorial supplement: K.\ Suominen, \emph{Localization of sheaves and Cousin complexes}, Acta Math.\ 131 (1973), 27--41, Theorem~3.9 and its following remark; \href{https://doi.org/10.1007/BF02392035}{doi:10.1007/BF02392035}. This result postdates the original book.) The specialization-stable filtration $Z^\bullet$ on the locally noetherian scheme $X$ is admissible and is strictly exhaustive because $d$ is bounded below. Pointwise dualizing complexes satisfy
\[
\mathcal H^i_{Z^p/Z^{p+1}}(R^\bullet)=0\qquad(i\ne p)
\]
by \bookref{motif:4.1.F}{IV,~1.F} and their local dualizing property. The Cousin equivalence therefore gives the natural isomorphism to $QE(R^\bullet)$ without an upper cohomology bound. Parts a) and b) identify its restrictions with the two categories in c).

For dualizing complexes, the corresponding target consists of residual complexes with the same $d$ and bounded terms. The printed target includes unbounded examples: on $X=\coprod_{n\geq0}\operatorname{Spec}k$, put $k$ on component $n$ in degree $n$. This residual complex has no finite global injective-dimension bound. The author's second remark is retained as historical commentary; the later fixed-filtration Cousin equivalence removes the bounded-cohomology obstruction in that setting.

\begin{exercise*}
Show that there is a uniqueness theorem analogous to \bookref{thm:5.3.1}{V,~3.1} for residual complexes, i.e., two residual complexes can differ only by shifting degrees and tensoring with an invertible sheaf. The touchy point is to show that if $K^{\bullet}, K^{\prime\bullet}$ are residual complexes, then $\sheafHom^{\bullet}(K^{\bullet}, K^{\prime\bullet})$ is a complex with coherent cohomology!\footnote{Editorial note: On a connected $X$, the shift is a single integer. On a general $X$, it is a locally constant integer-valued function, and the resulting complex must retain a global lower bound. The exercise is not supplied with a proof here.}

\end{exercise*}

We now give two technical results which will be used in the following sections.

\begin{booklemma}{1.2}\label{lem:6.1.2}
Let $K^\bullet$ and $K^{\prime\bullet}$ be residual complexes on $X$. Then to give an isomorphism $\psi\colon K^\bullet\to K^{\prime\bullet}$ is equivalent to giving, for each $x\in X$, an isomorphism
\[
\psi_x\colon Q(K_x^\bullet)\longrightarrow Q(K_x^{\prime\bullet})
\]
in $\DD_{\mathrm{c}}^+(\Spec\sO_{X,x})$, such that whenever $x\rightsquigarrow y$ is a specialization, then $\psi_x$ is obtained from $\psi_y$ by localization.
\end{booklemma}

\begin{proof}
Clearly $\psi$ gives rise to a system $(\psi_x)_{x\in X}$ as described. Conversely, given the isomorphisms $\psi_x$, we first note that by Proposition~\ref{prop:6.1.1}(c), $\psi_x$ comes from a unique isomorphism
\[
\overline{\psi}_x\colon K_x^\bullet\longrightarrow K_x^{\prime\bullet}
\]
of the actual residual complexes, and these $\overline{\psi}_x$ are compatible with localization. Already we deduce that the codimension functions $d$ and $d'$ associated to $K^\bullet$ and $K^{\prime\bullet}$ are the same. And since for $d(x)=d(y)$ and $x\neq y$ there are no non-zero maps of $J(x)$ into $J(y)$, our system of isomorphisms $\overline{\psi}_x$ gives rise to, and is determined by, a collection of isomorphisms
\[
\varphi_x\colon I(x)\longrightarrow I'(x)
\]
for each $x\in X$, where $I(x)$ (resp.\ $I'(x)$) is the (unique) subsheaf of $K^\bullet$ (resp.\ $K^{\prime\bullet}$) isomorphic to $J(x)$, such that for an immediate specialization $x\rightsquigarrow y$, $\varphi_x$ and $\varphi_y$ are compatible with the boundary maps of the complexes $K^\bullet$ and $K^{\prime\bullet}$. But to give such a collection of isomorphisms $(\varphi_x)$ is to give an isomorphism $\psi\colon K^\bullet\to K^{\prime\bullet}$, so we are done.\footnote{Editorial note: For the use of Proposition 1.1(c), first observe that the given local derived isomorphisms force equality of the codimension functions, by the degrees of the nonzero local Ext groups. Apply reconstruction with that common filtration; the subsequent chain-level isomorphisms are compatible with localization.}
\end{proof}

\begin{booklemma}[Glueing Lemma]{1.3}\label{lem:6.1.3}
Let $(U_\nu)$ be an open cover of $X$, and let $K_\nu^\bullet$ be a residual complex on $U_\nu$, for each $\nu$. Suppose furthermore that for each pair of indices $\mu,\nu$ we are given an isomorphism
\[
\alpha_{\mu\nu}\colon K_{\nu}^{\bullet}\longrightarrow K_{\mu}^{\bullet}
\]
of the restrictions of these complexes to $U_{\mu\nu}=U_{\mu}\cap U_{\nu}$, such that for each triple $\mu,\nu,\lambda$,
\[
\alpha_{\mu\nu}\alpha_{\nu\lambda}\alpha_{\lambda\mu}=1
\]
on $U_{\mu\nu\lambda}$, with $\alpha_{\nu\nu}=1$ and $\alpha_{\nu\mu}=\alpha_{\mu\nu}^{-1}$. Suppose finally that the lower bound of the complexes $K_{\nu}$ can be chosen uniformly for all $\nu$. Then there exists a residual complex, unique up to unique compatible isomorphism, $K^{\bullet}$ on $X$, together with isomorphisms
\[
\beta_{\nu}\colon K^{\bullet}|_{U_{\nu}}\longrightarrow K_{\nu}^{\bullet}
\]
for each $\nu$, which are compatible with $\alpha$ on $U_{\mu\nu}$.
\end{booklemma}

\begin{proof}
This is the usual situation for glueing, so it is clear that we can glue the complexes $K_{\nu}^{\bullet}$ into a global complex $K^{\bullet}$ on $X$, which is bounded below. It is also clear that $K^{\bullet}$ is a complex of quasi-coherent injective $\sO_X$-modules \bookref{lem:2.7.16}{II,~7.16}, and has coherent cohomology. Finally since the $J(x)$ are constant sheaves on irreducible subsets of $X$, we can glue together local isomorphisms to obtain a global isomorphism
\[
\bigoplus_{p\in\mathbb{Z}} K^{p} \cong \bigoplus_{x\in X} J(x).
\]
Hence $K^\bullet$ is a residual complex.
\end{proof}

\par\medskip\noindent\textbf{Remark.}
This lemma has the interesting consequence that one can glue dualizing complexes given locally in $\DD_{\mathrm{c}}^{+}(U_{\nu})$ into a global dualizing complex in $\DD_{\mathrm{c}}^{+}(X)$, whereas one cannot glue arbitrary objects of the derived category given locally (cf.\ remarks at end of \bookref{sec:3.8}{III,~\S8}).\footnote{Editorial note: The lemma permits gluing dualizing complexes supplied with compatible overlap isomorphisms satisfying the cocycle condition: pass to their residual complexes and glue the resulting actual complexes. The uniform lower bound in the lemma is essential. It yields a residual complex and hence a pointwise dualizing object of $\DD_{\mathrm{c}}^+(X)$. To conclude that this object is dualizing, one also needs a common upper bound for the injective dimensions of the local dualizing complexes. Their local Ext vanishing then gives global finite injective dimension by \bookref{prop:2.7.20}{II,~7.20}. These bounds are automatic for a finite cover, in particular after taking a finite subcover when $X$ is quasi-compact. (Editorial qualification of the original gluing claim: on $X=\coprod_{n\geq0}\operatorname{Spec}k$, the local complexes $k[n]$ have no uniform lower bound, while the local complexes $k[-n]$ glue to a bounded-below pointwise dualizing complex with no finite global injective-dimension bound. Thus neither bound can be omitted for arbitrary infinite covers.) This special gluing procedure does not assert descent for arbitrary objects of the derived category (cf.\ \bookref{sec:3.8}{III,~\S8}).}
\par\medskip

\par\endgroup

\section{Functorial properties}\label{sec:6.2}
\begingroup
\setlength{\emergencystretch}{2em}

In this section we define a notion of $f^{!}$ for residual complexes in case $f$ is a finite or smooth morphism. To avoid confusion, we introduce new, temporary notation $f^{y}$ and $f^{z}$ for these functors. The letters $y,z$ are arbitrary symbols; they have nothing to do with schemes $Y,Z$, or with points thereof. We give four isomorphisms I--IV between these functors, and seven compatibilities (i)--(vii) which will be referred to in the next two sections.

If $f\colon X\to Y$ is a \emph{finite} morphism of locally noetherian preschemes, we define a functor
\[
f^{y}\colon \operatorname{Res}(Y)\longrightarrow \operatorname{Res}(X)
\]
on the category of residual complexes by
\[
f^{y}(K^{\bullet}) = Ef^{\flat} Q(K^{\bullet}).
\]
Note by Proposition~\ref{prop:6.1.1}(a) that $Q(K^\bullet)$ is a pointwise dualizing complex; by \bookref{prop:5.2.4}{V,~2.4} $f^{\flat}Q(K^\bullet)$ is pointwise dualizing on $X$, and by Proposition~\ref{prop:6.1.1}(b), $Ef^{\flat}Q(K^\bullet)$ is a residual complex on $X$.

If $f\colon X\to Y$ and $g\colon Y\to Z$ are two finite morphisms of locally noetherian preschemes, we define an isomorphism
\[
(gf)^{y} \xrightarrow{\sim} f^{y}g^{y} \tag{I}\label{eq:6.2.I}
\]
of functors from $\operatorname{Res}(Z)$ to $\operatorname{Res}(X)$ as follows. By Lemma~\ref{lem:6.1.2} it will be enough to define this isomorphism in the scheme $\Spec(\sO_{X,x})$ for each $x\in X$, in a manner compatible with localization. A fortiori, it is enough to define the isomorphism after making a base extension $\Spec(\sO_{Z,z})\to Z$, for each point $z\in Z$, and thus we reduce to the case where $Z$ (and hence also $X$ and $Y$) are noetherian of finite Krull dimension. In that case pointwise dualizing complexes are dualizing, and we have a functorial isomorphism $\varphi\colon 1\to QE$ on the category of dualizing complexes on $Y$ (Proposition~\ref{prop:6.1.1}(c)). Now for a residual complex $K^\bullet$ on $Z$ we define our isomorphism as follows:
Using the composition isomorphism \bookref{prop:3.6.2}{III,~6.2}, we set
\[
\begin{aligned}
(gf)^yK^\bullet
 &= E(gf)^\flat QK^\bullet
  \cong Ef^\flat g^\flat QK^\bullet \\
 &\xrightarrow[\varphi]{\sim} Ef^\flat QE g^\flat QK^\bullet
  = f^y g^y(K^\bullet).
\end{aligned}
\]

We will use this same technique of reduction to the case of finite Krull dimension without explicit mention below. It enables us to carry over isomorphisms defined for dualizing complexes to residual complexes.

\phantomsection\label{compat:6.2.i}
For a composition of three finite morphisms, there is the usual commutative diagram (referred to as (i)) of the isomorphisms \eqref{eq:6.2.I}.

For a \emph{smooth} morphism $f\colon X\to Y$ of locally noetherian preschemes we define a functor
\[
f^{z}\colon \operatorname{Res}(Y)\longrightarrow \operatorname{Res}(X)
\]
by
\[
f^{z}(K^\bullet) = Ef^{\sharp}Q(K^\bullet).
\]
This takes residual complexes into residual complexes by virtue of Proposition~\ref{prop:6.1.1}(a),(b), and \bookref{thm:5.8.3}{V,~8.3}.

For a composition of two smooth morphisms $f,g$, we define an isomorphism
\[
(gf)^{z} \xrightarrow{\sim} f^{z}g^{z} \tag{II}\label{eq:6.2.II}
\]
using the above reduction to the case of finite Krull dimension, and carrying over the isomorphism \bookref{prop:3.2.2}{III,~2.2}.

\phantomsection\label{compat:6.2.ii}
For a composition of three smooth morphisms, there is a compatibility (referred to as (ii)) of the isomorphisms \eqref{eq:6.2.II}, which follows from the compatibility of \bookref{prop:3.2.2}{III,~2.2}.

There are two other isomorphisms, expressing compatibilities between the functors $f^{y}$ and $g^{z}$: the Cartesian square, and the residue isomorphism. For the first, we suppose given a Cartesian diagram as shown:
\[
\begin{tikzcd}
	W \arrow[r, "k"] \arrow[d, "h"'] & Y \arrow[d, "g"] \\
	X \arrow[r, "f"'] & Z
\end{tikzcd}
\]
(i.e., $W = X\times_Z Y$) with $f$ (and hence $k$) a finite morphism, and $g$ (and hence $h$) a smooth morphism. In that case there is an isomorphism
\[
h^{z}f^{y} \xrightarrow{\sim} k^{y}g^{z} \tag{III}\label{eq:6.2.III}
\]
obtained as above from the isomorphism of \bookref{cor:3.6.4}{III,~6.4}.

If we have another Cartesian diagram as shown:
\[
\begin{tikzcd}
	Q \arrow[r, "p"] \arrow[d, "q"'] & W \arrow[r, "k"] \arrow[d, "h"'] & Y \arrow[d, "g"] \\
	P \arrow[r, "r"'] & X \arrow[r, "f"'] & Z
\end{tikzcd}
\]
with $r$ also a finite morphism, then there is a commutative diagram of isomorphisms
\[
\begin{tikzcd}[row sep=2.5em, column sep=4em]
	q^z(fr)^y
	\arrow[r, "\text{III}"]
	\arrow[d, "\text{I}"']
	& (kp)^y g^z
	\arrow[dr, "\text{I}"] \\
	q^z r^y f^y
	\arrow[r, "\text{III}"]
	& p^y h^z f^y
	\arrow[r, "\text{III}"]
	& p^y k^y g^z
\end{tikzcd}
\tag{iii}\label{eq:6.2.iii}
\]
\phantomsection\label{compat:6.2.iv}
Also if we have another Cartesian diagram doubling the smooth side of the square,
\[
\begin{tikzcd}
	Q \arrow[r, "p"] \arrow[d, "q"'] & P \arrow[d, "r"] \\
	W \arrow[r, "k"] \arrow[d, "h"'] & Y \arrow[d, "g"] \\
	X \arrow[r, "f"'] & Z
\end{tikzcd}
\]
i.e., with $r$ smooth, then there is a similar commutative diagram (iv) involving the isomorphisms \eqref{eq:6.2.II} and (III). These two compatibilities follow from \bookref{cor:3.6.4}{III,~6.4}.

For the \emph{residue isomorphism}, we suppose that we have a finite morphism $f$ which is factored into $pi$, with $i$ a closed immersion, and $p$ smooth.
\[
\begin{tikzcd}
	& P \arrow[dd, "p"] \\
	X \arrow[ur, "i"] \arrow[dr, "f"'] & \\
	& Y
\end{tikzcd}
\]
Then there is an isomorphism
\[
f^{y} \xrightarrow{\sim} i^{y}p^{z} \tag{IV}\label{eq:6.2.IV}
\]
obtained as above from that of \bookref{prop:3.8.2}{III,~8.2}.

If we have a diagram with $f,g$ finite, $i,k$ closed immersions, and $q,p$ smooth,
\[
\begin{tikzcd}
X\arrow[rr, "i"] \arrow[d, "f"]&&A\arrow[dl, "q"]\\
Y\arrow[r, "k"] \arrow[d, "g"]&P\arrow[dl, "p"]\\
Z
\end{tikzcd}
\]
then there is a commutative diagram
\[
\begin{tikzcd}
(gf)^{y} \ar[r,"\mathrm{IV}"] \ar[d,"\mathrm{I}"'] & i^{y}(pq)^{z} \ar[r,"\mathrm{II}"] & i^{y}q^{z}p^{z} \\
f^{y}g^{y} \ar[r,"\mathrm{IV}"] & f^{y}k^{y}p^{z} \ar[r,"\mathrm{I}"] & (kf)^{y}p^{z} \ar[u,"\mathrm{IV}"']
\end{tikzcd}
\tag{v}\label{eq:6.2.v}
\]
This follows from \bookref{prop:3.8.6}{III,~8.6(b)} applied to the triples $(i,q,p)$ and $(f,k,p)$.

There are also two commutative diagrams expressing compatibility between the square isomorphism \eqref{eq:6.2.III} and the residue isomorphism \eqref{eq:6.2.IV}.

For the first, suppose one has a diagram with $Q = Y\times_{Z} P$, $i$ finite, $j,k,\ell$ closed immersions, and $f,g$ smooth.
\[
\begin{tikzcd}
& Q\arrow[dd, "g"]\arrow[r, "l"]&P\arrow[dd, "f"]\\
X\arrow[ru, "k"]\arrow[rd, "i"]\\
&Y\arrow[r, "j"]&Z
\end{tikzcd}
\]


Then there is a commutative diagram (vi)
\[
\begin{tikzcd}[row sep=2.5em, column sep=4em]
	(ji)^y
	\arrow[r, "\text{IV}"]
	\arrow[d]
	& (\ell k)^y f^z
	\arrow[dr, "\text{I}"] \\
	i^y j^y
	\arrow[r, "\text{IV}"]
	& k^y g^z j^y
	\arrow[r, "\text{III}"]
	& k^y \ell^y f^z
\end{tikzcd}
\tag{vi}\label{eq:6.2.vi}
\]
which follows from \bookref{prop:3.8.6}{III,~8.6(c)}.

For the second, suppose there is a Cartesian diagram (with $W = X\times_{Y} Z$; $Q = P\times_{Y} Z$) with $f,g$ finite, $i,j$ closed immersions, and $p,q,u,v,w$ smooth.
\[
\begin{tikzcd}
	& Q \arrow[dd, "v", pos=0.3] \arrow[rd, "q"] & \\
	W \arrow[ur, "j"] \arrow[rr, "g"', pos=0.6] \arrow[dd, "u"'] & & Z \arrow[dd, "w"] \\
	& P\arrow[rd, "p"] & \\
	X \arrow[ur, "i"] \arrow[rr, "f"'] & & Y
\end{tikzcd}
\]
Then there is a commutative diagram
\[
\begin{tikzcd}
u^{z}f^{y} \ar[r,"\mathrm{IV}"] \ar[d,"\mathrm{III}"'] & u^{z}i^{y}p^{z} \ar[r,"\mathrm{III}"] & j^{y}v^{z}p^{z}  \ar[dr,"\mathrm{II}"']\\
g^{y}w^{z} \ar[r,"\mathrm{IV}"] & j^{y}q^{z}w^{z} \ar[r,"\mathrm{II}"] & j^{y}(wq)^{z} \ar[r,"="]& j^{y}(pv)^{z}
\end{tikzcd}
\tag{vii}\label{eq:6.2.vii}
\]
This follows from \bookref{prop:3.8.6}{III,~8.6(a)} modified as in \bookref{cor:3.6.4}{III,~6.4}.

\par\endgroup

\section{$f^{!}$ for residual complexes}\label{sec:6.3}
\begingroup
\setlength{\emergencystretch}{2em}

In this section we construct the functor $f^{!}$ for residual complexes, which we call $f^{\Delta}$ to avoid confusion. We refer back to the previous section for the notations $f^{y}$, $f^{z}$, the isomorphisms \eqref{eq:6.2.I}--\eqref{eq:6.2.IV}, and the compatibilities \hyperref[compat:6.2.i]{(i)}--\eqref{eq:6.2.vii}.

We will work in the category of \emph{locally noetherian preschemes}, and we will consider only morphisms which are of \emph{finite type}, and such that the \emph{dimensions of the fibres are bounded}. It will be understood in the following that these conditions hold for all schemes and morphisms considered.

We are now in a position to state our theorem.

\begin{booktheorem}{3.1}\label{thm:6.3.1}
There exists a theory of variance consisting of the data a)--d) below subject to the conditions VAR 1 -- VAR 5. Furthermore, this theory is unique in the sense that given a second collection of such data a')--d') there is a unique isomorphism of the functors a) and a') compatible with the isomorphisms b)--d) and b')--d').

\begin{enumerate}
\item[a)] For every morphism $f\colon X\to Y$, a functor
\[
f^{\Delta}\colon \operatorname{Res}(Y)\longrightarrow \operatorname{Res}(X)
\]
on the category of residual complexes.
\item[b)] For every pair of consecutive morphisms $f\colon X\to Y$ and $g\colon Y\to Z$, an isomorphism
\[
c_{f,g}\colon (gf)^{\Delta} \xrightarrow{\sim} f^{\Delta}g^{\Delta}.
\]
\item[c)] For every finite morphism $f$, an isomorphism
\[
\psi_f\colon f^{\Delta} \xrightarrow{\sim} f^{y}.
\]
\item[d)] For every smooth morphism $g$, an isomorphism
\[
\phi_g\colon g^{\Delta} \xrightarrow{\sim} g^{z}.
\]
\end{enumerate}

\medskip\phantomsection\label{var:6.3.1}\noindent\textbf{VAR 1).} For any $f$, $c_{\id,f}=c_{f,\id}=\id$, and if $f,g,h$ are three consecutive morphisms, then there is a commutative diagram
\[
\begin{tikzcd}
(hgf)^{\Delta} \ar[r,"c_{f,hg}"] \ar[d,"c_{gf,h}"'] & f^{\Delta}(hg)^{\Delta} \ar[d,"c_{g,h}"] \\
(gf)^{\Delta}h^{\Delta} \ar[r,"c_{f,g}"] & f^{\Delta}g^{\Delta}h^{\Delta}
\end{tikzcd}
\]
(In other words, the $f^{\Delta}$ and $c_{f,g}$ define a ``cat\'egorie cliv\'ee normalis\'ee'' in the terminology of \cite[\nopp 1960--61, expos\'e~VI]{SGA}).

\medskip\phantomsection\label{var:6.3.2}\noindent\textbf{VAR 2).} If $f,g$ are consecutive finite morphisms, then $c_{f,g}$ is compatible, via $\psi_f$ and $\psi_g$, with the usual isomorphism \eqref{eq:6.2.I} above.

\medskip\phantomsection\label{var:6.3.3}\noindent\textbf{VAR 3).} If $f,g$ are consecutive smooth morphisms, then $c_{f,g}$ is compatible, via $\phi_f$ and $\phi_g$, with the usual isomorphism \eqref{eq:6.2.II} above.

\medskip\phantomsection\label{var:6.3.4}\noindent\textbf{VAR 4).} Given a Cartesian diagram
\[
\begin{tikzcd}
	W \arrow[r, "k"] \arrow[d, "h"'] & Y \arrow[d, "g"] \\
	X \arrow[r, "f"'] & Z
\end{tikzcd}
\]
with $f,k$ finite and $g,h$ smooth, the isomorphisms b) are compatible via c), d) with the isomorphism \eqref{eq:6.2.III} of a square above, i.e., there is a commutative diagram
\[
\begin{tikzcd}[row sep=3em, column sep=4em]
	(fh)^\Delta = (gk)^\Delta
	\arrow[rr, "c_{k,g}"]
	\arrow[d, "c_{h,f}"']
	&& k^\Delta g^\Delta
	\arrow[d, "{\psi_k, \phi_g}"] \\
	h^\Delta f^\Delta
	\arrow[r, "{\phi_h, \psi_f}"]
	& h^z f^y
	\arrow[r, "\text{III}"]
	& k^y g^z
\end{tikzcd}
\]

\medskip\phantomsection\label{var:6.3.5}\noindent\textbf{VAR 5).} Given a diagram
\[
\begin{tikzcd}
X \ar[rr,"f"] \ar[dr,"i"'] & & Y \\
& P \ar[ur,"p"']
\end{tikzcd}
\]
with $i$ a closed immersion, $f$ finite, and $p$ smooth, we have a commutative diagram
\[
\begin{tikzcd}
f^{\Delta}=(pi)^{\Delta} \ar[r,"c_{i,p}"] \ar[d,"\psi_f"'] & i^{\Delta}p^{\Delta}\ar[d,"{\psi_i,\phi_p}"] \\
f^{y} \ar[r,"\mathrm{IV}"] & i^{y}p^{z}
\end{tikzcd}
\]
i.e., $c_{i,p}$ is compatible with the residue isomorphism.
\end{booktheorem}

The proof of this theorem requires drawing a great many diagrams and checking their commutativity. We will therefore carry out in detail only a few of these verifications, by way of example, and will leave the others to the reader, marking them with the symbol (!) which indicates that he has some work to do at that point. The proof per se will follow after some definitions and lemmas.

\begin{definition*}
Let $f\colon X\to Y$ be a fixed morphism, and let $U\subseteq X$ be an open set. We define a \emph{chart} on $U$ to be the following collection of data:
\begin{enumerate}
\item[1)] A functor $f^{a}\colon \operatorname{Res}(Y)\to \operatorname{Res}(U)$.
\item[2)] A factorization $f|_{U} = pi$ where $i$ is a closed immersion into a scheme $P$ smooth over $Y$.
\[
\begin{tikzcd}
U \ar[rr,"f|_{U}"] \ar[dr,"i"'] & & Y \\
& P \ar[ur,"p"']
\end{tikzcd}
\]
\item[3)] An isomorphism
\[
\gamma_{i,p}\colon f^{a} \xrightarrow{\sim} i^{y}p^{z}.
\]
\end{enumerate}
\end{definition*}

\begin{definition*}
If $f^{b}$, $f|_{U} = qj$, $q\colon Q\to Y$, and $\gamma'_{j,q}$ is a second chart on the same open set $U$, a \emph{permissible isomorphism} between the two is an isomorphism
\[
\alpha\colon f^{a} \longrightarrow f^{b}
\]
of functors, such that for every commutative diagram
\[
\begin{tikzcd}
	& R \arrow[ld, "r"'] \arrow[rd, "s"] & \\
	P \arrow[rd, "p"'] & U \arrow[l, "i"'] \arrow[r, "j"] \arrow[u, "k"] \arrow[d, "f"'] & Q \arrow[ld, "q"] \\
	& Y &
\end{tikzcd}
\]
with $k$ a closed immersion and $r,s$ smooth, there is a commutative diagram of isomorphisms
\[
\begin{tikzcd}
f^{a} \ar[r,"\gamma_{i,p}"] \ar[d,"\alpha"'] & i^{y}p^{z} \ar[r,"\mathrm{IV}"] & k^{y}r^{z}p^{z}\ar[dr,"II"] \\
f^{b} \ar[r,"\gamma'_{j,q}"] & j^{y}q^{z} \ar[r,"\mathrm{IV}"] & k^{y}s^{z}q^{z} \ar[r,"\mathrm{II}"] & k^{y}(qs)^{z} = k^{y}(pr)^{z}
\end{tikzcd}
\]
\end{definition*}

\par\endgroup

\begingroup
\setlength{\emergencystretch}{2em}
\begin{booklemma}{3.2}\label{lem:6.3.2}
Given two charts on an open set $U$, there exists a unique permissible isomorphism between them.
\end{booklemma}

\begin{proof}
The uniqueness is clear, because one can take $R = P\times_{Y} Q$ and $k$ the diagonal map $i\times j$. Then the isomorphism $\alpha$ is determined by the condition of the definition.

For the existence, let $f^{a},f^{b}$, etc., be two charts as above, and let $(R,k,r,s)$ and $(S,\ell,\cdots)$ be two diagrams as above. We must show that the isomorphisms $\alpha_R$ and $\alpha_S$ defined by the condition above are the same. By considering the third diagram $(R\times_Y S, k\times\ell,\cdots)$ and comparing each to this one, we reduce to the case where $S$ dominates $R$. In other words, we have a commutative diagram with $k,\ell$ closed immersions and $r,s,t$ smooth morphisms.
\[
\begin{tikzcd}
	& S \arrow[rrd, "t"] &&& \\
	&&& R \arrow[rd, "s"] \arrow[ldd, "r"',pos=0.7] & \\
	U \arrow[ruu, "\ell"]
	\arrow[rrru, "k"]
	\arrow[rrrr, "j"]
	\arrow[rrd, "i"']
	\arrow[ddrr, "f"']
	&&&& Q \arrow[ddll, "q"] \\
	&& P \arrow[d, "p"] && \\
	&& Y &&
\end{tikzcd}
\]
We must show that the following diagram is commutative:
\[
\begin{tikzcd}[row sep=1.2cm, column sep=0.45cm]
	f^a \arrow[r, "\gamma_{i,p}"] \arrow[d, "\mathrm{id}"] &
	i^y p^z \arrow[r, "\mathrm{IV}"] &
	k^y r^z p^z \arrow[r, "\mathrm{II}"] &
	k^y(pr)^z = k^y(qs)^z &
	k^y s^z q^z \arrow[l, "\mathrm{II}"'] &
	j^y q^z \arrow[l, "\mathrm{IV}"'] &
	f^b \arrow[l, "\gamma'_{j,q}"'] \arrow[d, "\mathrm{id}"] \\
	f^a \arrow[r, "\gamma_{i,p}"] &
	i^y p^z \arrow[r, "\mathrm{IV}"] &
	\ell^y(rt)^z p^z \arrow[r, "\mathrm{II}"] &
	\ell^y(prt)^z = \ell^y(qst)^z &
	\ell^y(st)^z q^z \arrow[l, "\mathrm{II}"'] &
	j^y q^z \arrow[l, "\mathrm{IV}"'] &
	f^b \arrow[l, "\gamma'_{j,q}"']
\end{tikzcd}
\]
We will check this one completely by filling in more arrows to make little commutative diagrams. At the left hand end of the diagram we fill in as follows:
\[
\begin{tikzcd}
f^{a} \ar[r,"\gamma_{i,p}"] \ar[dd,"\id"'] & i^{y}p^{z} \ar[r,"\mathrm{IV}"] \ar[dd,"\id"'] & k^{y}r^{z}p^{z} \ar[d,"\mathrm{IV}"] \\
&&\ell^{y}t^{z}r^{z}p^{z}\ar[d,"\mathrm{II}"]\\
f^{a} \ar[r,"\gamma_{i,p}"] & i^{y}p^{z} \ar[r,"\mathrm{IV}"] & \ell^{y}(rt)^{z}p^{z}
\end{tikzcd}
\]
The left-hand square is identically commutative. The right-hand one is commutative by virtue of a special case of compatibility \eqref{eq:6.2.v} above, where $X=Y$.

On the right-hand side of our long diagram we fill in two analogous commutative squares. This leaves in the middle the following diagram

	\[
	\begin{tikzcd}[row sep=1.4cm, column sep=1.4cm]
		% Row 1
		k^y r^z p^z \arrow[r, "\mathrm{II}"] \arrow[d, "\mathrm{IV}"] &
		k^y(pr)^z \arrow[r, phantom, "="] \arrow[d, "\mathrm{IV}"] &
		k^y(qs)^z \arrow[d, "\mathrm{IV}"] &
		k^y s^z q^z \arrow[l, "\mathrm{II}"'] \arrow[d, "\mathrm{IV}"] \\
		% Row 2
		\ell^y t^z r^z p^z \arrow[r, "\mathrm{II}"] &
		\ell^y t^z(pr)^z \arrow[r, phantom, "="] &
		\ell^y t^z(qs)^z &
		\ell^y t^z s^z q^z \arrow[l, "\mathrm{II}"'] \\
		% Row 3
		\ell^y(rt)^z p^z \arrow[r, "\mathrm{II}"] \arrow[u, "\mathrm{II}"'] &
		\ell^y(prt)^z \arrow[r, phantom, "="] \arrow[u, "\mathrm{II}"'] &
		\ell^y(qst)^z \arrow[u, "\mathrm{II}"'] &
		\ell^y(st)^z q^z \arrow[l, "\mathrm{II}"'] \arrow[u, "\mathrm{II}"']
	\end{tikzcd}
	\]

The middle squares are obviously commutative; the upper left and upper right are commutative because the isomorphism IV operates between $U$ and $R$, and II operates between $R$ and $Y$, so that the order doesn't matter; and the lower left and lower right are commutative by \hyperref[compat:6.2.ii]{(ii)} above.

\end{proof}

\begin{booklemma}{3.3}\label{lem:6.3.3}
The composition of permissible isomorphisms is permissible. The inverse of a permissible isomorphism is permissible.
\end{booklemma}

\begin{proof}
The inverse is obvious. For the composition, let $f^{a}$ and $f^{b}$ be two charts as above, and let $f^{c}$, $\ell$, $r$, $R$ be a third. Using $P\times_Y Q\times_Y R$ to construct the unique permissible isomorphisms of $f^{a}$ with $f^{b}$, $f^{b}$ with $f^{c}$, and $f^{a}$ with $f^{c}$, we see that the composition of the first two is the third.
\end{proof}

\par\endgroup

\begin{proof}[Proof of Theorem~\ref{thm:6.3.1}]
\begingroup
\setlength{\emergencystretch}{2em}
First we prove the existence of the theory of variance.

\medskip\phantomsection\label{construction:6.3.a}\noindent a) \textit{Construction of the functor $f^{\Delta}$}. Let $f\colon X\to Y$ be a morphism. Let $U = (U_{\nu})$ be a cover of $X$ by open sets with charts on them. Note that charts exist locally: If $x\in X$ is a point, let $V$ be a noetherian affine neighborhood of $f(x)$ in $Y$, and let $U$ be an affine neighborhood of $x$ in $f^{-1}(V)$, which is of finite type over $V$. Then $U$ admits a closed immersion $i$ into an affine space $\mathbf{A}_{V}^{n}$ for suitable $n$. Let $P=\mathbf{A}_{V}^{n}$, let $p\colon P\to Y$ be the natural projection, let $f^{a}=i^{y}p^{z}$, and $\gamma_{i,p}=\id$. This gives a chart on $U$ which is a neighborhood of the given point $x$.

If $f^{a}$, $i$, $P$, $p$, $\gamma_{i,p}$ is a chart on an open set $U$, and if $U'\subseteq U$ is a smaller open set, we define the notion of a \emph{restriction} (not unique) of the chart to $U'$, as follows: let $P'\subseteq P$ be an open set whose intersection with $i(U)$ is $i(U')$. Then we take $f^{a}|_{U'}$, $i|_{U'}$, $P'$, $p|_{P'}$, $\gamma_{i,p}|_{U'}$ as the restriction of the chart. Note that a restriction of permissible isomorphisms is permissible.

For each pair of indices $\mu,\nu$, choose restrictions of the charts on $U_{\nu}$ and $U_{\mu}$ to $U_{\mu\nu}=U_{\nu}\cap U_{\mu}$, and let $\alpha_{\mu\nu}$ be the unique permissible isomorphism between them. One sees immediately that the isomorphism of functors
\[
\alpha_{\mu\nu}\colon f^{\nu}\longrightarrow f^{\mu}
\]
thus defined on $U_{\mu\nu}$ is independent of the restrictions of the charts chosen. Furthermore it follows from the lemmas that on a triple intersection $U_{\mu\nu\lambda}$, these isomorphisms are compatible:
\[
\alpha_{\mu\nu}\alpha_{\nu\lambda}\alpha_{\lambda\mu} = \id.
\]
The lower bounds needed for gluing are uniform. Indeed, if the input residual complex has lower bound $a$ and the fibre dimensions of $f$ are at most $N$, the finite and smooth codimension formulas (in particular \bookref{cor:5.8.4}{V,~8.4}) show that the local output complexes have lower bound $a-N$.

Thus, since we are dealing with functors on residual complexes, we can use Lemma~\ref{lem:6.1.3}, and glue the functors $f^{\nu}$ via the isomorphisms $\alpha_{\mu\nu}$ to obtain a functor
\[
f^{\Delta}\colon \operatorname{Res}(Y)\longrightarrow \operatorname{Res}(X)
\]
(together with isomorphisms $\beta_{\nu}\colon f^{\Delta}|_{U_{\nu}}\to f^{\nu}$ for each $\nu$ compatible with the $\alpha_{\mu\nu}$) which is the one we want.

\medskip\phantomsection\label{construction:6.3.b}\noindent b) \textit{Construction of the isomorphisms $c_{f,g}$}. Let $f\colon X\to Y$ and $g\colon Y\to Z$ be morphisms. It will be sufficient to construct $c_{f,g}$ locally, chart by chart, and show that it is compatible with the permissible isomorphisms of change of chart. For then we can glue. Thus we may assume that $f$ and $g$ are embeddable in smooth morphisms, and we may even assume that $f$ is embeddable in an affine space over $Y$, since that is possible locally. Choose embeddings
$i: X \longrightarrow \mathbb{A}^n_Y$ and $k: Y \longrightarrow P$, smooth
over $Z$, and note that $\mathbb{A}^n_Y = \mathbb{A}^n_P \times_P Y$, so
that we have a commutative diagram as
shown, with $i, j, k$ closed immersions, $p, q, r$ smooth, and the upper right square Cartesian.
\[
\begin{tikzcd}[row sep=1.2cm, column sep=1cm]
	X \arrow[r, "i"] \arrow[dr, "f"'] & \mathbb{A}^n_Y \arrow[rr, "j"] \arrow[d, "p"'] & & \mathbb{A}^n_P \arrow[d, "q"] \\
	& Y \arrow[rr, "k"] \arrow[dr, "g"'] & & P \arrow[dl, "r"] \\
	& & Z &
\end{tikzcd}
\]

Define $c_{f,g}$ (depending on the charts chosen, and taking $ji: X \longrightarrow \mathbb{A}^n_P$ as a chart for $gf$) as follows:

\[
\begin{tikzcd}
	c_{f,g}: (gf)^a \arrow[r, "\gamma_{ji,rq}"] & (ji)^y(rq)^z \arrow[r,"\mathrm{I,II}"] & i^y j^y q^z r^z & i^y p^z k^y r^z \arrow[l, "\mathrm{III}"'] & f^b g^c. \arrow[l, "\gamma_{i,p} \gamma_{k,r}"']
\end{tikzcd}
\]
Now we must show that $c_{f,g}$ is compatible with the permissible isomorphisms of change of chart for $f$ and $g$. It is sufficient to vary one at a time. Furthermore, if $k'\colon Y\to Q$ is another chart for $g$, by considering the third chart $k'' = k\times k'\colon Y\to P\times_Z Q$, we reduce to the case where $Q$ dominates $P$. So we have diagram such as the one shown,

\[\begin{tikzcd}[row sep=2.5em, column sep=2.5em]
	& & & \mathbb{A}^n_Q \arrow[rd, "v"] & \\
	& & \mathbb{A}^n_Y \arrow[ru, "j'"] \arrow[rr, "j" near start] \arrow[dd, "p"'] & & \mathbb{A}^n_P \arrow[dd, "q"] \\
	X \arrow[rru, "i"] \arrow[rrd, "f"'] & & & Q \arrow[rd, "u"] & \\
	& & Y \arrow[ru, "k'"'] \arrow[rr, "k"'] \arrow[rd, "g"'] & & P \arrow[ld, "r"] \\
	& & & Z &
	\arrow[from=1-4, to=3-4, "q'" near start, crossing over]
\end{tikzcd}\]
and we must show that the two $c_{f,g}$ are compatible, i.e., that the following diagram is commutative:
For legibility, write
\[
\begin{gathered}
\Gamma=\gamma_{ji,rq},\qquad \Gamma'=\gamma_{j'i,rqv},\qquad
\theta=\gamma_{i,p},\\
\eta=\gamma_{k,r},\qquad \eta'=\gamma_{k',ru}.
\end{gathered}
\]
\[\begin{tikzcd}[column sep=1em]
	(gf)^a \arrow[r, "\Gamma"] \arrow[d, "\alpha_{gf}"'] & (ji)^y(rq)^z \arrow[r, "\mathrm{I,II}"] & i^y j^y q^z r^z & i^y p^z k^y r^z \arrow[l, "\mathrm{III}"'] & i^y p^z g^c \arrow[l, "\eta"'] & f^b g^c \arrow[l, "\theta"'] \arrow[d, "\alpha_{g}"] \\
	(gf)^d \arrow[r, "\Gamma'"] & (j'i)^y(ruq')^z \arrow[r, "\mathrm{I,II}"] & i^y j'^y q'^z (ru)^z & i^y p^z k'^y (ru)^z \arrow[l, "\mathrm{III}"'] & i^y p^z g^e \arrow[l, "\eta'"'] & f^b g^e \arrow[l, "\theta"']
\end{tikzcd}\]

One may verify this (!) using the definition of a permissible isomorphism, and the compatibilities \hyperref[compat:6.2.ii]{(ii)}, \eqref{eq:6.2.v}, and \eqref{eq:6.2.vii} above.

Similarly we must consider a change of the chart on $X$ to another, say $i'\colon X\to\mathbf{A}_Y^{m}$, and as before we may assume that $i'$ dominates $i$. This involves checking another analogous commutative diagram of isomorphisms, which we leave to the reader (!).

\par\endgroup

\begingroup
\setlength{\emergencystretch}{2em}
\medskip\phantomsection\label{construction:6.3.c}\noindent c) \textit{Construction of the isomorphism $\psi_f$ for a finite morphism $f$}. Let $f\colon X\to Y$ be a finite morphism.
\[\begin{tikzcd}
	X\arrow[rr, "i"]\arrow[rd, "f"]&&P\arrow[ld, "p"]\\
	&Y
\end{tikzcd}\]
As before, it will be sufficient to construct the isomorphism $\psi_f$ locally on charts, provided our definition is compatible with permissible isomorphisms of charts. So let $i\colon X\to P$ be a chart for $f$, and define
\[
\psi_f\colon f^{a} \xrightarrow{\gamma_{i,p}} i^{y}p^{z} \xleftarrow{\mathrm{IV}} f^{y}.
\]
If $j\colon X\to Q$ is another chart, one may assume as usual that $j$ dominates $i$, and one checks (!) using compatibility \eqref{eq:6.2.v} above that $\psi_f$ is compatible.

\medskip\phantomsection\label{construction:6.3.d}\noindent d) \textit{For a smooth morphism $g\colon X\to Y$} we must construct the isomorphism $\phi_g$. Take $g$ as its own chart, and take $\phi_g = \gamma$. There is no choice involved, hence nothing to check.

Having constructed the data a)--d), we must verify the conditions VAR 1 -- VAR 5.

\medskip\noindent\textbf{VAR 1).} Given three morphisms $f,g,h$, we use a diagram such as the one shown to calculate $c_{f,g}$, etc., on the charts. The
diagram of the condition becomes a large diagram whose commutativity one checks (!) using \eqref{eq:6.2.iii} and \hyperref[compat:6.2.iv]{(iv)} above.
\[
\begin{tikzcd}[column sep=1.5em, row sep=1.5em]
	X \arrow[rr] \arrow[rd, "f"'] && \mathbb{A}^m_Y \arrow[rr] \arrow[ld] && \mathbb{A}^{m+n}_Z \arrow[rr] \arrow[ld] && \mathbb{A}^{m+n}_P \arrow[ld] \\
    &Y \arrow[rr] \arrow[rd, "g"'] && \mathbb{A}^n_Z \arrow[rr] \arrow[ld] && \mathbb{A}^n_P \arrow[ld] \\
	& &  Z \arrow[rr] \arrow[rd, "h"'] && P \arrow[ld] \\
	& & & W
\end{tikzcd}
\]
\medskip\noindent\textbf{VAR 2).} If $f$ and $g$ are consecutive finite morphisms, one checks (!) that $c_{f,g}$ is compatible with \eqref{eq:6.2.I} via the isomorphisms $\psi_f$ and $\psi_g$, by using \eqref{eq:6.2.v} and \eqref{eq:6.2.vi} above.

\medskip\noindent\textbf{VAR 3).} Trivial. One has only to observe that the isomorphism \eqref{eq:6.2.III} is the identity if one side of the square is the identity.

\medskip\noindent\textbf{VAR 4) and VAR 5)} follow (!) from the definitions and a few more commutative diagrams.

This completes the proof of existence of the theory of variance, and we now show its uniqueness. For that purpose, suppose that $\{f^{\Delta}, c^{\Delta}_{f,g}, \psi^{\Delta}_{f}, \phi^{\Delta}_{g}\}$ and $\{f^{X}, c^{X}_{f,g}, \psi^{X}_{f}, \phi^{X}_{g}\}$ are two sets of data a)--d), each satisfying the conditions VAR 1 -- VAR 5 (which we will call VAR $1^{\Delta}$, VAR $1^{X}$, etc., to be precise). We will construct an isomorphism $\delta\colon f^{\Delta}\to f^{X}$, compatible with the data b)--d) and b')--d'), and we will observe that $\delta$ is unique.

Let $f\colon X\to Y$ be a morphism of finite type, and choose locally an embedding $i\colon X\to P$ into a scheme smooth over $Y$.

\[\begin{tikzcd}
	X\arrow[rr, "i"]\arrow[rd, "f"]&&P\arrow[ld, "p"]\\
	&Y
\end{tikzcd}\]

Define $\delta$ by
\[
\delta\colon f^{\Delta} \xrightarrow{c^{\Delta}_{i,p}} i^{\Delta} p^{\Delta} \xrightarrow{\psi^{\Delta}_{i}, \phi^{\Delta}_{p}} i^{y} p^{z} \xleftarrow{\psi^{X}_{i}, \phi^{X}_{p}} i^{X} p^{X} \xleftarrow{c^{X}_{i,p}} f^{X}.
\]
(Note that in order to be compatible with the $c_{i,p}$, $\psi_i$, $\phi_p$, we must choose $\delta$ this way, which proves the uniqueness of $\delta$.)

\par\endgroup

To see that $\delta$ is independent of the embedding chosen, let $j\colon X\to Q$ be another one. Replacing $Q$ by $P\times_Y Q$ as usual we reduce to the case $Q$ dominates $P$, and have a diagram such as the one shown.

\[
\begin{tikzcd}
& Q \ar[rd,"q"] \\
X \ar[ru,"j"] \ar[rr,"i"] \ar[rd,"f"'] && P \ar[ld,"p"] \\
&Y
\end{tikzcd}
\]

We must check that the perimeter of the resulting diagram is commutative.

Write, for $T=\Delta,X$,
\[
A_T=(\psi_i^T,\phi_p^T),\qquad
B_T=(\psi_j^T,\phi_q^T,\phi_p^T),\qquad
C_T=(\psi_j^T,\phi_{pq}^T).
\]
\[
\begin{tikzcd}[column sep=1.2em,row sep=4em]
f^\Delta \ar[r,"c^\Delta_{i,p}"] \ar[dd,"\id"'] &
i^\Delta p^\Delta \ar[r,"{A_{\Delta}}"] \ar[d,"c^\Delta_{j,q}"'] &
i^y p^z \ar[d,"\mathrm{IV}"] &
i^X p^X \ar[l,"{A_{X}}"'] \ar[d,"c^X_{j,q}"] &
f^X \ar[l,"c^X_{i,p}"'] \ar[dd,"\id"] \\
& j^\Delta q^\Delta p^\Delta \ar[r,"{B_{\Delta}}"] &
j^y q^z p^z &
j^X q^X p^X \ar[l,"{B_{X}}"'] & \\
f^\Delta \ar[r,"c^\Delta_{j,pq}"'] &
j^\Delta(pq)^\Delta \ar[u,"c^\Delta_{q,p}"] \ar[r,"{C_{\Delta}}"'] &
j^y(pq)^z \ar[u,"\mathrm{II}"'] &
j^X(pq)^X \ar[u,"c^X_{q,p}"'] \ar[l,"{C_{X}}"] &
f^X \ar[l,"c^X_{j,pq}"]
\end{tikzcd}
\]

The left inner regions commute by VAR $1^{\Delta}$, VAR $3^{\Delta}$, and VAR $5^{\Delta}$, together with the residue and smooth-composition isomorphisms IV and II. The right regions use VAR $1^{X}$, VAR $3^{X}$, and VAR $5^{X}$ in the same way.

Thus $\delta$ is well-defined locally. Now cover $X$ by open sets $U$ for which this is possible, and glue together the local isomorphisms
\[
\delta_{\nu}\colon (f j_{\nu})^{\Delta} \xrightarrow{\sim} (f j_{\nu})^{X}
\]
where $j_{\nu}\colon U_{\nu}\to X$ is the open immersion. Since $j_{\nu}^{\Delta}$ is isomorphic via $\phi_{j_{\nu}}^{\Delta}$ with $j_{\nu}^{z}$ which is the restriction, we can glue the $\delta_{\nu}$ once we have checked (!) that the isomorphisms are compatible with restriction. Thus
\[
\delta\colon f^{\Delta} \xrightarrow{\sim} f^{X}
\]
is defined.

Now we must check that $\delta$ is compatible with the isomorphisms $c_{f,g}$, $\psi_f$, $\phi_g$. This can be done locally using the conditions VAR 1 -- VAR 5, and we leave the details to the reader (!).
\end{proof}

\begin{bookproposition}{3.4}\label{prop:6.3.4}
Let $f\colon X\to Y$ be a morphism (with the conventions above) and let $K^\bullet$ be a residual complex on $Y$. Let $d$ denote the codimension function on $X$ (resp.\ $Y$) associated with the pointwise dualizing complex $f^{\Delta}K^\bullet$ (resp.\ $K^\bullet$). Then for each $x\in X$, if $y=f(x)$, we have
\[
d(x) = d(y) - \operatorname{tr.d.} k(x)/k(y).
\]
\end{bookproposition}

\begin{proof}
The question is local, and compatible with composition of morphisms. Thus we reduce to the case $f$ finite (trivial) or $f$ smooth (which is \bookref{cor:5.8.4}{V,~8.4}).
\end{proof}

\begin{bookcorollary}{3.5}\label{cor:6.3.5}
Let $f\colon X\to Y$ be a morphism (with the conventions above: $f$ is of finite type, and the dimensions of its fibres are bounded). Then if $Y$ admits a residual complex (resp.\ dualizing complex) so does $X$.
\end{bookcorollary}

\begin{proof}
The existence of a residual complex on $X$, given one on $Y$, follows from Theorem~\ref{thm:6.3.1}. If $Y$ admits a dualizing complex $R^\bullet$, then the associated codimension function $d$ on $Y$ is bounded. By Proposition~\ref{prop:6.3.4}, the codimension function associated to $Qf^{\Delta}E(R^\bullet)$ is also bounded. But $Qf^{\Delta}E(R^\bullet)$ is pointwise dualizing, hence dualizing (cf.\ proof of \bookref{prop:5.8.2}{V,~8.2}).
\end{proof}

\section{Trace for residual complexes}\label{sec:6.4}

In this section we define the trace map for residual complexes. For a morphism $f\colon X\to Y$ (with the conventions of the previous section) it is a map of graded sheaves
\[
\Tr_f\colon f_*f^{\Delta}K^\bullet \longrightarrow K^\bullet
\]
where $K^\bullet$ is a residual complex on $Y$. It will be a morphism of complexes only if $f$ is proper (see the Residue Theorem in the next chapter, which generalizes the classical theorem that the sum of the residues of a differential on a curve is zero).

First we define the trace map for a finite morphism, by carrying over the map of \bookref{prop:3.6.5}{III,~6.5} to residual complexes. Let $f\colon X\to Y$ be a finite morphism. We will denote by $f'$ the functor
\[
\overline{f}^{*}\sheafHom_{\sO_Y}(f_*\sO_X,\cdot)
\]
(using the notation of \bookref{sec:3.6}{III,~\S6}) so that $f^{\flat} = \RR f'$.

\begin{booklemma}{4.1}\label{lem:6.4.1}
Let $f\colon X\to Y$ be a finite morphism, and let $K^\bullet$ be a residual complex on $Y$. Then $f'(K^\bullet)$ is a residual complex on $X$, and $f_*f'(K^\bullet)$ is a Cousin complex on $Y$, with respect to the filtration $Z^\bullet$ associated with $Q(K^\bullet)$ (cf.\ \bookref{sec:4.3}{IV,~\S3}).
\end{booklemma}

\begin{proof}
It is clear that $f'$ takes quasi-coherent injectives on $Y$ to quasi-coherent injectives on $X$, and it is also clear that $f'(K^\bullet)$ has coherent cohomology, so we have only to check that it is isomorphic to a sum of sheaves of the form $J(x)$. Indeed, it will be sufficient to show, for each $y\in Y$, that
\[
f'(J(y)) \cong \bigoplus_{f(x)=y} J(x).
\]
Letting $A$ be the local ring of $y$ on $Y$, and $B$ the semilocal ring of the points $x$ of $X$ lying over $y$, we have a homomorphism from the local ring $A$ to the semilocal ring $B$, where $B$ is a finite $A$-module. If $k$ is the residue field of $A$, and $I$ is an injective hull of $k$ over $A$, we must show that $J = \Hom_A(B,I)$ is a direct sum $\bigoplus_{i=1}^{r} I_i$ where $k_1,\dots,k_r$ are the residue fields of $B$, and $I_i$ is an injective hull of $k_i$ over $B$. For any $B$-module of finite type $M$, we have isomorphisms
\[
\Hom_{B}(M,J) = \Hom_{B}(M, \Hom_{A}(B,I)) \cong \Hom_{A}(M,I).
\]
Now $J$ is injective, and has support at the closed points of $\Spec B$. Hence it is a direct sum of some number of copies of the injectives $I_i$. To find out how many, we have only to calculate
\[
\Hom_{B}(k_i,J) \cong \Hom_{A}(k_i,I) \cong k_i
\]
which is true since $k_i$ is an $A$-module of finite length, and $I$ is dualizing. Hence each $I_i$ occurs just once.

Now $f_*f'(K^\bullet)$ is a direct sum, for each $y\in Y$, of
\[
\bigoplus_{f(x)=y} f_* J(x)
\]
which is a constant sheaf spread out on $\overline{\{y\}}$, and which occurs in degree $p = d(y)$, where $d$ is the associated codimension function. Hence $f_*f'(K^\bullet)$ is a Cousin complex with respect to the filtration $Z^\bullet$ associated with $Q(K^\bullet)$.

\end{proof}

\begingroup
\setlength{\emergencystretch}{2em}
Now we are in a position to define the trace map for the finite morphism $f$, which we will call $\rho_f$ to avoid confusion:
\[
\rho_{f}\colon f_{*}f^{y}(K^{\bullet}) \longrightarrow K^{\bullet}.
\]
Let $K^\bullet$ be a residual complex on $Y$. Since it is an injective complex, the natural map
\[
\xi_{f}\colon Qf'(K^\bullet) \longrightarrow f^{\flat}Q(K^\bullet)
\]
of the functor $f'$ into its derived functor $f^{\flat}$ is an isomorphism in $\DD^{+}(X)$ \bookref{thm:1.5.1}{I,~5.1}. Hence also
\[
E\xi_{f}\colon f'(K^\bullet) = EQf'(K^\bullet) \longrightarrow Ef^{\flat}Q(K^\bullet) = f^{y}(K^\bullet)
\]
is an isomorphism of residual complexes (recall $EQ \cong 1$ for any Cousin complex).

Consider the map
\[
Qf_{*}f'(K^{\bullet}) \xrightarrow{\xi_{f_{*}}} \RR f_{*}Qf'(K^{\bullet}) \xrightarrow{\xi_{f'}} \RR f_{*}f^{\flat}Q(K^{\bullet}) \xrightarrow{\operatorname{Trf}_{f}} Q(K^{\bullet}),
\]
where the last map, $\operatorname{Trf}_f$, is the trace of \bookref{prop:3.6.5}{III,~6.5}. By Lemma~\ref{lem:6.4.1}, $f_*f'(K^\bullet)$ is a Cousin complex with respect to the filtration $Z^\bullet$, and $K^\bullet$ is an injective Cousin complex, so by \bookref{lem:4.3.2}{IV,~3.2}, this map in the derived category is represented by a unique map of complexes
\[
f_*f'(K^\bullet) \longrightarrow K^\bullet.
\]
Composing with the inverse of the isomorphism $E\xi_f$ above, we obtain the desired map
\[
\rho_{f}\colon f_{*}f^{y}(K^{\bullet}) \longrightarrow K^{\bullet}.
\]

If $f\colon X\to Y$, and $g\colon Y\to Z$ are two finite morphisms, then there is a commutative diagram
\[
\begin{tikzcd}
(gf)_*(gf)^{y} \ar[r,"\rho_{gf}"] \ar[d,"\mathrm{I}"'] & 1 \\
g_*f_* f^{y}g^{y} \ar[r,"\rho_f"] & g_*g^{y} \ar[u,"\rho_g"']
\end{tikzcd}
\tag{viii}\label{eq:6.4.viii}
\]
which follows from \bookref{prop:3.6.6}{III,~6.6} and the functoriality of the above construction.

\par\medskip\noindent\textbf{Remark.}
We have given the above definition in some detail to establish its relation with the trace map defined for the derived category in Chapter III. One could of course define $\rho_f$ much more quickly by the usual ``evaluation at one'' map, without passing to the derived category, but we need the functorial properties below.
\par\medskip

\begin{booktheorem}{4.2}\label{thm:6.4.2}
As above, we work in the category of locally noetherian preschemes, and morphisms of finite type such that the dimension of the fibres is bounded. There exists a unique theory of trace, consisting of the data a) below, subject to the conditions TRA 1 and TRA 2.

\begin{enumerate}
\item[a)] For each morphism $f\colon X\to Y$, a morphism $\Tr_{f}\colon f_{*}f^{\Delta}\to 1$ of functors from $\operatorname{Res}(Y)$ to the category of graded sheaves of $\sO_Y$-modules (where $1$ denotes the forgetful functor: consider a residual complex $K_Y^{\bullet}$ simply as a graded $\sO_Y$-module).
\end{enumerate}

\medskip\phantomsection\label{tra:6.4.1}\noindent\textbf{TRA 1).} If $f$ and $g$ are two consecutive morphisms of finite type, then there is a commutative diagram
\[
\begin{tikzcd}
(gf)_*(gf)^{\Delta} \ar[r,"\Tr_{gf}"] \ar[d,"c_{f,g}"'] & 1 \\
g_*f_* f^{\Delta}g^{\Delta} \ar[r,"\Tr_f"] & g_*g^{\Delta} \ar[u,"\Tr_g"']
\end{tikzcd}
\]

\medskip\phantomsection\label{tra:6.4.2}\noindent\textbf{TRA 2).} If $f$ is a finite morphism, then $\Tr_f$ is the one we already know, i.e., $\Tr_f = \rho_f\circ f_*(\psi_f)$, using the notation of Theorem~\ref{thm:6.3.1}.
\end{booktheorem}

\par\endgroup

\begingroup
\setlength{\emergencystretch}{2em}
\begin{booklemma}{4.3}\label{lem:6.4.3}
\cite[\nopp IV,~\S?]{EGA} Let $f\colon X\to Y$ be a morphism of finite type, let $x\in X$ be a point which is closed in its fibre, and let $Z$ be the closure of $x$ in $X$, with the reduced induced structure. Then there is an open neighborhood $V$ of $y=f(x)$ such that $Z\cap f^{-1}(V)$ is finite over $V$.
\end{booklemma}

\begin{proof}
Replacing $X$ by $Z$ and $Y$ by $\overline{f(Z)}$ with the reduced induced structure, we reduce to the case where $X$ and $Y$ are integral schemes, $x$ is the generic point of $X$, and $f$ is dominant. Since furthermore the question is local on $X$ and $Y$, we may assume $X$ and $Y$ are affine, and thus we reduce to the following statement in algebra:\footnote{Editorial note: First restrict $Y$ to an affine neighborhood of $y$. Replacing $X$ by $Z$ and $Y$ by the reduced closure of $y$, we may assume that $X$ and $Y$ are integral, $Y$ is affine noetherian with generic point $y$, and $f$ is dominant. The generic fibre has the single point $x$, since $x$ was closed in its original fibre. An open neighborhood of $y$ in this reduced closure is obtained by intersecting it with an open neighborhood in the original target.\par
Choose an affine open $U\subseteq X$ containing $x$, and let $W=X\setminus U$ with its reduced closed subscheme structure. The image $f(W)$ is constructible by Chevalley's theorem, \cite[\nopp IV,~1.8.5]{EGA}. Since it does not contain $y$, its closure is a proper closed subset of the integral scheme $Y$. After shrinking $Y$ to a nonempty principal open disjoint from that closure, the restricted source is contained in $U$ and is a principal open subset of $U$, hence affine. (Editorial completion: the original cites only ``EGA IV \S?'' and says the question is local on $X$ and $Y$. The shrinking step above justifies the affine reduction. The pinpoint citation identifies its constructibility input, not a claimed verbatim occurrence of the whole lemma.) We have therefore reduced to the following statement in algebra:}

Let $A\to B$ be an inclusion of integral domains, with $B$ a finitely generated $A$-algebra. Let $K$ be the quotient field of $A$, and assume that $B\otimes_A K$ is a finite extension field of $K$. Then there exists an element $f\neq0$ in $A$ such that $B\otimes_A A_f$ is a finite $A_f$-module.

To prove this statement, let $b_1,\dots,b_n$ be a finite set of elements in $B$, with $b_1=1$, such that
\begin{enumerate}
\item[1)] the $b_i$ generate $B$ as an $A$-algebra, and
\item[2)] the elements $b_i\otimes 1$ generate $B\otimes K$ as a $K$-vector space.
\end{enumerate}
Then for each $i,j$, $b_i b_j\otimes 1 \in B\otimes K$, so we can write
\[
b_i b_j \otimes 1 = \sum_{k=1}^{n} b_k \otimes \lambda_{ijk},\qquad \lambda_{ijk}\in K.
\]
Let $f$ be a common denominator for all the $\lambda_{ijk}$, $1\leq i,j,k\leq n$. Then the $\lambda_{ijk}$ are all in $A_f$, so the $A_f$-module $B_0$ generated by the $b_i\otimes 1$ is in fact a ring. But $B_0$ contains a set of generators for $B\otimes A_f$ as an $A_f$-algebra, so $B_0 = B\otimes A_f$, which shows that the latter is a finitely generated $A_f$-module.
\end{proof}

\begin{booklemma}{4.4}\label{lem:6.4.4}
Let $f\colon X\to Y$ be a morphism, and let $K^\bullet$ be a residual complex on $Y$. Then there is a unique map
\[
\Tr_f\colon f_*f^{\Delta}K^{\bullet} \longrightarrow K^{\bullet}
\]
of graded sheaves on $Y$, such that whenever $U$ is an open subset of $Y$, and $W$ is a closed subscheme of $f^{-1}(U)$, finite over $U$
\[
\begin{tikzcd}
W \ar[rr,"i"] \ar[dr,"g"'] && f^{-1}(U) \ar[dl,"f"] \\
& U
\end{tikzcd}
\]
 we require the following compatibility. Write $L^\bullet=f^\Delta K^\bullet$ and $N^\bullet=g^\Delta K^\bullet$. The required diagram is
\[
\begin{tikzcd}[column sep=1.2em,row sep=3em]
g_*g^y K^\bullet \ar[drr,"\rho_g"'] &
g_*N^\bullet \ar[l,"\psi_g"'] \ar[r,"c_{i,f}"] &
f_*i_*i^\Delta L^\bullet \ar[r,"\psi_i"] &
f_*i_*i^y L^\bullet \ar[r,"\rho_i"] &
f_*L^\bullet \ar[dll,"\Tr_f"] \\
&& K^\bullet &&
\end{tikzcd}
\]
(Here we denote by $f$ the restriction of $f$ to $f^{-1}(U)$, by $K^\bullet$ the restriction of $K^\bullet$ to $U$, and so forth.)
\end{booklemma}\par\endgroup


\begin{proof}
\begingroup
\setlength{\emergencystretch}{2em}
Let $d$ denote the codimension function on $Y$ associated with the residual complex $K^{\bullet}$, and let it also denote the codimension function on $X$ associated with the residual complex $f^{\Delta}K^{\bullet}$ (the one we mean will always be clear from the context). Then according to the definition of the residual complexes we have isomorphisms
\[
K^{p} \cong \bigoplus_{d(y)=p} J(y)
\]
and
\[
(f^{\Delta}K^{\bullet})^{p} \cong \bigoplus_{d(x)=p} J(x).
\]
We will denote by $J(y)$ (resp.\ $J(x)$) the unique subsheaf of $K^{p}$ (resp.\ $f^{\Delta}(K^{\bullet})^{p}$) which is isomorphic to an injective hull of $k(y)$ (resp.\ $k(x)$), so as to make these isomorphisms canonical.

I claim that to give a map of graded $\sO_Y$-modules
\[
\Tr_f\colon f_*f^{\Delta}K^{\bullet} \longrightarrow K^{\bullet}
\]
is equivalent to giving, for each $x\in X$ which is closed in its fibre, a map
\[
\Tr_{f,x}\colon f_{*}J(x) \longrightarrow J(y)
\]
where $y=f(x)$. Indeed, $x$ is closed in its fibre if and only if $\operatorname{tr.d.} k(x)/k(y)=0$, i.e., if and only if $d(x)=d(y)$, by Proposition~\ref{prop:6.3.4} above. Hence $\Tr_f$ certainly gives rise to a collection of maps $\Tr_{f,x}$ as above. On the other hand, if $y'\in Y$ is a point not equal to $y$, with $d(y')=d(x)$, then there is no non-zero map of $f_*J(x)$ into $J(y')$, because $d(y')\leq d(y)$ by Proposition~\ref{prop:6.3.4}, and so $y'\notin\overline{\{y\}}$. (Note that $f_*J(x)$ has support in $\overline{\{y\}}$.) Thus $\Tr_f$ is determined by the maps $\Tr_{f,x}$ above, and these can be given arbitrarily.

We will now construct maps $\Tr_{f,x}$ as above for each point $x\in X$ which is closed in its fibre.

Given an $x\in X$ closed in its fibre, choose by Lemma~\ref{lem:6.4.3} an open neighborhood $V$ of $y=f(x)$ such that $Z = \overline{\{x\}}\cap f^{-1}(V)$ is finite over $V$.
\[
\begin{tikzcd}
Z \ar[rr,"i"] \ar[dr,"g"'] && f^{-1}(V) \ar[dl,"f"] \\
&V
\end{tikzcd}
\]
Let $\sJ$ be the ideal of $Z$ with the reduced induced structure, and let $Z_n$, $n=1,2,\dots$, be the subscheme of $X$ defined by $\sJ^{n}$.
\[
\begin{tikzcd}
& Z_{n'} \ar[rd,"i_{n'}"] \ar[dd,"g_{n'}",near start] & \\
Z_n \ar[ru,"j_{nn'}"] \ar[rr,"i_n",pos=0.8] \ar[rd,"g_n"'] && f^{-1}(V) \ar[ld,"f"] \\
&V&
\end{tikzcd}
\]
For $n<n'$ we have a closed immersion $j_{nn'}\colon Z_n\to Z_{n'}$, and one can verify (!) using \hyperref[var:6.3.1]{VAR 1}, \hyperref[var:6.3.2]{VAR 2}, and \eqref{eq:6.4.viii} above, that the following diagrams are commutative.
For this diagram put $F^\bullet=f^\Delta K^\bullet$ and
\[
\begin{gathered}
A_m=f_*i_{m*}i_m^y F^\bullet,\qquad B_m=g_{m*}g_m^y K^\bullet\quad(m=n,n'),\\
C=f_*i_{n'*}j_{nn'*}j_{nn'}^y i_{n'}^y F^\bullet,\qquad T=f_*F^\bullet.
\end{gathered}
\]
Let $\theta_m\colon B_m\xrightarrow{\sim}A_m$ be the isomorphism obtained from $\psi_{g_m}^{-1}$, $c_{i_m,f}$, and $\psi_{i_m}$, with the indicated pushforwards. Then the comparison is
\[
\begin{tikzcd}[column sep=3em,row sep=3em]
& B_n \ar[r,"\theta_n"] \ar[dl,"\rho_{g_n}"'] & A_n \ar[d,"\mathrm{I}","\sim"'] \ar[dr,"\rho_{i_n}"] & \\
K^\bullet && C \ar[d,"\rho_{j_{nn'}}"] & T \\
& B_{n'} \ar[r,"\theta_{n'}"'] \ar[ul,"\rho_{g_{n'}}"] & A_{n'} \ar[ur,"\rho_{i_{n'}}"'] &
\end{tikzcd}
\]

\par\endgroup

One can also write down a diagram with three indices $n,n',n''$ showing that the sheaves $A_n=f_*i_{n*}i_n^y f^\Delta(K^\bullet)$ and the transition maps induced by $\rho_{j_{nn'}}$ and I form a direct system. The maps $\rho_{i_n}$ embed this system into $f_*f^\Delta(K^\bullet)$, and the compatible maps
\[
\rho_{g_n}\circ\theta_n^{-1}\colon A_n\longrightarrow K^\bullet
\]
map it to $K^\bullet$. Thus we may pass to the limit and obtain a map from a subsheaf of $f_*f^\Delta(K^\bullet)$ to $K^\bullet$.

On the component $f_*J(x)$, the $n$th term, via the inclusion $\rho_{i_n}$, is
\[
f_*\sheafHom_{\sO_X}(\sO_X/\sJ^n,J(x)).
\]
Its limit is $f_*J(x)$ itself, since $J(x)$ has support on $Z$. Hence the resulting map is defined on all of $f_*J(x)$, and we define it to be $\Tr_{f,x}$.


To complete the proof of the lemma, we must check various things. First, that $\Tr_{f,x}$ is well-defined, i.e., does not depend on the choice of the open set $V$. That is clear, since $\Tr_{f,x}$ depends only on the stalks at $x$ and $y$, and all of our constructions are compatible with localization.

Second, we must check the property of the lemma. So let $W$ and $U$ be as in the statement. It will be sufficient to check the diagram on the component of $x$ for each $x\in W$. Let $Z$ be the closure of $x$, with the reduced induced structure (which will be finite over $U$). Let $J$ be the ideal of $Z$ in $W$, and for each $n=1,2,\ldots$, let $Z_n'$ be the subscheme of $W$ defined by $J^{n}$. Let $Z_n$ be as above. Then there are closed immersions as shown. We leave to the reader (!) to write down a huge commutative diagram of $Z_n$'s which in the limit gives the diagram of the lemma on the component of $x$.
\[\begin{tikzcd}
	W\ar[r]&f^{-1}(U)\\
	Z'_n\ar[r]\ar[u]&Z_n\ar[u]
\end{tikzcd}\]
Third we must check the uniqueness, but this is clear from the construction.
\end{proof}

\begin{proof}[Proof of Theorem~\ref{thm:6.4.2}]
For a morphism $f\colon X\to Y$, we define $\Tr_f$ to be the map given by Lemma~\ref{lem:6.4.4}. It is clearly functorial in $K^\bullet$. In case $f$ is a finite morphism, we can take $U=Y$ and $W=X$, so that $\Tr_f = \rho_f\circ f_*(\psi_f)$, which proves condition \hyperref[tra:6.4.2]{TRA 2}.
\[\begin{tikzcd}
	W\ar[r,"i"]\ar[d]&X\ar[d,"f"]\\
	V\ar[r,"j"]\ar[rd]&Y\ar[d,"g"]\\
	&Z\\
\end{tikzcd}\]
To prove \hyperref[tra:6.4.1]{TRA 1}, let $f\colon X\to Y$ and $g\colon Y\to Z$ be two morphisms. It is sufficient to prove the condition of TRA 1 for a single residual complex $K^\bullet$ on $Z$ and for a single $x\in X$ which is closed in its fibre over $Z$. The question is local, so we may assume that $W = \overline{\{x\}}$ is finite over $Z$, and that $V = \overline{\{y\}}$, where $y=f(x)$, is finite over $Z$. Finally it is sufficient to check the commutativity for a given element $a\in\Gamma(J(x))$. We choose a subscheme structure $W_n$ on $W$ with $n$ large enough so that $a$ is in the image of $i^{y}J(x)$. Now one can check (!) the required commutativity using the property of the lemma, and \eqref{eq:6.4.viii}, \hyperref[var:6.3.1]{VAR 1}, and \hyperref[var:6.3.2]{VAR 2} above.

The uniqueness of $\Tr_f$ is clear, as it was in the lemma, since TRA 1 and TRA 2 imply the condition of the lemma.

\end{proof}

\section{Behavior with respect to certain base changes}\label{sec:6.5}
\begingroup
\setlength{\emergencystretch}{2em}

In this section we show that the functor $f^{\Delta}$ and the morphism $\Tr_{f}$ are compatible with certain base changes which take residual complexes into residual complexes.

\begin{definition*}
A morphism $f\colon X\to Y$ of locally noetherian preschemes is \emph{residually stable} if
\begin{enumerate}
\item[a)] $f$ is flat
\item[b)] $f$ is integral \cite[\nopp II,~6.1.1]{EGA}, and
\item[c)] the fibres of $f$ are Gorenstein preschemes \bookref{sec:5.9}{V,~\S9}.
\end{enumerate}
\end{definition*}

\begin{example*}
\begin{enumerate}
\item An open immersion is residually stable.\footnote{Editorial note: An open-and-closed immersion is residually stable: it is flat as an open immersion, integral as a closed immersion, and its nonempty fibres are spectra of fields. (Editorial correction: the original says ``an open immersion''. Under the integral definition above, an open immersion is residually stable if and only if its image is also closed, since integral morphisms are closed. For example, $\Spec\mathbb Z[1/2]\to\Spec\mathbb Z$ is open but is not integral.)
\par For an arbitrary open immersion $j\colon U\hookrightarrow Y$, restriction still takes residual complexes to residual complexes. Indeed, $j^*J(y)$ is $J_U(y)$ if $y\in U$, and zero otherwise; the defining decomposition, coherence and lower bound therefore restrict to $U$. This restriction fact does not require $j$ to satisfy the integral definition above.}
\item A composition of residually stable morphisms is residually stable (use \bookref{prop:5.9.6}{V,~9.6}).
\item If $X$ and $Y$ are the spectra of fields $k\subseteq K$, then $f$ is residually stable if and only if $K$ is algebraic over $k$.
\item If $f\colon X\to Y$ is residually stable, then every point $x\in X$ is closed in its fibre, since $f$ is integral. Thus the fibres are zero-dimensional Gorenstein preschemes.
\end{enumerate}
\end{example*}

\begin{bookproposition}{5.1}\label{prop:6.5.1}
If $f\colon X\to Y$ is a residually stable morphism, and if $u\colon Y'\to Y$ is a morphism of finite type, then $g\colon X' = X\times_Y Y'\to Y'$ is also residually stable.
\end{bookproposition}

\begin{proof}
Clearly $g$ is flat and integral. To show that the fibres of $g$ are Gorenstein, let $y'\in Y'$, let $y = u(y')$, and consider the map of fibres $v\colon X'_{y'}\to X_{y}$. Now $X_{y}$ is Gorenstein since $f$ is residually stable. The fibres of $v$ are tensor products of fields, one of which is finitely generated (namely $k(y')/k(y)$), hence Gorenstein \bookref{cor:5.9.5}{V,~9.5}. Hence $X'_{y'}$ is Gorenstein \bookref{prop:5.9.6}{V,~9.6}.
\end{proof}

\begin{booklemma}{5.2}\label{lem:6.5.2}
Let $A$ be a noetherian local ring with maximal ideal $\mathfrak{m}$, and let $I$ be an $A$-module. Then $I$ is an injective hull of the residue field $k$ of $A$ if and only if
\begin{enumerate}
\item[a)] $I$ has support at the closed point of $\Spec A$,
\item[b)] $\Hom_{A}(k,I)\cong k$, and
\item[c)] There is a sequence of ideals $\mathfrak{a}_1\supseteq\mathfrak{a}_2\supseteq\cdots$ of $A$ which form a base for the $\mathfrak{m}$-adic topology, such that for all $n$, $\operatorname{length}(A/\mathfrak{a}_n) = \operatorname{length}(\Hom_A(A/\mathfrak{a}_n, I))$.
\end{enumerate}
\end{booklemma}

\begin{proof}
If $I$ is an injective hull of $k$, then a), b), c) are immediate, since $\Hom_A(\cdot,I)$ is a dualizing functor for modules of finite length (Proposition~\ref{prop:5.5.1}). In fact, c) holds for any ideal which is primary for $\mathfrak{m}$.

Conversely, let $I$ be an $A$-module satisfying a), b), c), and let $J$ be an injective hull of $k$. By b), we can find an injection $k\subseteq I$, and since $J$ is injective, the natural map $\theta\colon k\to J$ extends to a map $\psi\colon I\to J$. I claim $\psi$ is injective. Indeed, given $y\in I$, $y\neq0$, we have $\mathfrak{m}^n y=0$ for sufficiently large $n$, by a). Let $n$ be the least such integer. Then there is an $x\in\mathfrak{m}^{n-1}$ for which $xy\neq0$. But $\mathfrak{m}(xy)=0$, so $xy\in k$. Hence $\psi(xy)=\theta(xy)\neq0$. But $\psi(xy)=x\cdot\psi(y)$, so $\psi(y)\neq0$, and $\psi$ is injective.

Now choose a sequence of ideals $\mathfrak{a}_1\supseteq\mathfrak{a}_2\supseteq\cdots$ as in c). Then we note that
\[
\psi(\Hom(A/\mathfrak{a}_n,I)) \subseteq \Hom(A/\mathfrak{a}_n,J),
\]
and both have the same length, namely the length of $A/\mathfrak{a}_n$. Therefore they are equal. But $I$ and $J$ are the union of the submodules of elements annihilated by $\mathfrak{a}_n$, so we see that $\psi$ is also surjective. Thus $I$ is isomorphic to $J$, and so is an injective hull of $k$.
\end{proof}

\par\endgroup

\begingroup
\setlength{\emergencystretch}{2em}
\begin{bookproposition}{5.3}\label{prop:6.5.3}
Let $f\colon X\to Y$ be a residually stable morphism, and let $K^\bullet$ be a residual complex on $Y$. Then $f^{*}(K^\bullet)$ is a residual complex on $X$.
\end{bookproposition}

\begin{proof}
Clearly $f^{*}(K^\bullet)$ is a complex of quasi-coherent sheaves with coherent cohomology. We have only to check that there is an isomorphism
\[
\bigoplus_{p} f^{*}(K^{p}) \cong \bigoplus_{x\in X} J(x).
\]
Thus we reduce to the following local statement: For each $x\in X$, let $y=f(x)$, and let $I$ be an injective hull of $k=k(y)$ over the local ring $A=\sO_{Y,y}$ of $y$. Let $B=\sO_{X,x}$ be the local ring of $x$. Then $I\otimes_A B$ is an injective hull of $k'=k(x)$ over $B$.

We apply the criteria of Lemma~\ref{lem:6.5.2}. Since $I$ is an injective hull of $k$ over $A$, we have
\begin{enumerate}
\item[a)] $I$ has support at the closed point of $\Spec A$
\item[b)] $\Hom_{A}(k,I)\cong k$
\item[c)] For each $n$, $\operatorname{length}(A/\mathfrak{m}^{n}) = \operatorname{length}(\Hom_A(A/\mathfrak{m}^{n}, I))$,
\end{enumerate}
where $\mathfrak{m}$ is the maximal ideal of $A$.

Now since $x$ is closed in its fibre, $I\otimes_A B$ has support at the closed point of $\Spec B$. Since $B$ is flat over $A$, we have
\[
\Hom_{B}(B/\mathfrak{m}B, I\otimes_A B) \cong B/\mathfrak{m}B.
\]
Therefore
\[
\Hom_{B}(k', I\otimes_A B) \cong \Hom_{B}(k', B/\mathfrak{m}B) \cong k',
\]
since $B/\mathfrak{m}B$ is an Artinian Gorenstein ring. Finally note that for each $n$,
\[
\operatorname{length}(B/\mathfrak{m}^{n}B) = \operatorname{length}(\Hom_B(B/\mathfrak{m}^{n}B, I\otimes_A B)).
\]
Thus the conditions of Lemma~\ref{lem:6.5.2} are fulfilled, and $I\otimes_A B$ is an injective hull of $k'$ over $B$.
\end{proof}

\begin{bookproposition}{5.4}\label{prop:6.5.4}
Let $Y$ be the spectrum of a local Artin ring $A$. Then one can find a residually stable morphism $u\colon Y'\to Y$ where $Y'$ is the spectrum of a local Artin ring $A'$ with algebraically closed residue field.
\end{bookproposition}

\begin{proof}
By \cite[\nopp $0_{\mathrm{III}}$,~10.3.1]{EGA} one can find a local Artin ring $A'$ and a local homomorphism $A\to A'$ such that $A'$ is flat over $A$, $\mathfrak{m}A' = \mathfrak{m}'$, and $A'/\mathfrak{m}'$ is an algebraic closure of $A/\mathfrak{m}$. Then $Y' = \Spec A'$ will do. Indeed, $Y'$ is flat over $Y$. It is integral since it is Artinian, and its residue field is algebraic over that of $Y$. The only fibre is a field, which is Gorenstein.
\end{proof}

\par\endgroup

\begingroup
\setlength{\emergencystretch}{2em}
Now we come to the behavior of $f^{\Delta}$ and $\Tr$ under residually stable base change. Let $f\colon X\to Y$ be a morphism of finite type, and let $u\colon Y'\to Y$ be a residually stable morphism. Let $X' = X\times_{Y} Y'$, and let $u', f'$ be as shown:
\[
\begin{tikzcd}
X' \ar[r,"u'"] \ar[d,"f'"'] & X \ar[d,"f"] \\
Y' \ar[r,"u"] & Y
\end{tikzcd}
\]

If $f$ is a finite morphism, there is a natural isomorphism
\[
f'^{y}u^{*} \xrightarrow{\sim} u'^{*}f^{y} \tag{V}\label{eq:6.5.V}
\]
derived from \bookref{prop:3.6.3}{III,~6.3} (cf.\ section~\hyperref[sec:6.2]{2} above).

On the other hand if $f$ is smooth, there is a natural isomorphism
\[
f'^{z}u^{*} \xrightarrow{\sim} u'^{*}f^{z} \tag{VI}\label{eq:6.5.VI}
\]
derived from \bookref{prop:3.2.1}{III,~2.1}.

\begin{booktheorem}{5.5}\label{thm:6.5.5}
For every morphism $f\colon X\to Y$ of finite type, and every residually stable morphism $u\colon Y'\to Y$ (using the notations above) there is an isomorphism
\[
d_{u,f}\colon f'^{\Delta}u^{*} \xrightarrow{\sim} u'^{*}f^{\Delta}
\]
such that
\begin{enumerate}
\item[1)] If $v\colon Y''\to Y'$ is another residually stable morphism, then
\[
d_{uv,f}=(v'^*d_{u,f})\circ(d_{v,f'}u^*),
\]
where $v'\colon X''\to X'$ is the induced base change.
\item[2)] If $f,g$ are two consecutive morphisms of finite type, then the isomorphisms $d$ are compatible with $c_{f,g}$ and $c_{f',g'}$.
\item[3)] If $f$ is a finite morphism, then $d_{u,f}$ is compatible with \eqref{eq:6.5.V} via the isomorphisms $\psi_f,\psi_{f'}$.
\item[4)] If $f$ is a smooth morphism, then $d_{u,f}$ is compatible with \eqref{eq:6.5.VI} via the isomorphisms $\phi_f,\phi_{f'}$.
\end{enumerate}
\end{booktheorem}

\begin{proof}[Suggestion of Proof (!)]
Show first that V and VI are compatible with the isomorphisms I, II, III and IV of section~\hyperref[sec:6.2]{2}. Define $d_{u,f}$ for $f$ finite or smooth using 3) and 4). Then define $d_{u,f}$ for arbitrary $f$ and check its properties, following the construction of $f^{\Delta}$ given in section~\hyperref[sec:6.3]{3}.
\end{proof}

\begin{booktheorem}{5.6}\label{thm:6.5.6}
Let $u$ and $f$ be as in Theorem~\ref{thm:6.5.5}. Then there is a commutative diagram
\[
\begin{tikzcd}
f'_* f'^{\Delta} u^{*} \ar[r,"\Tr_{f'}"] \ar[d,"d_{u,f}"'] & u^{*} \\
f'_* u'^{*} f^{\Delta} \ar[r,"\sim"] & u^{*} f_* f^{\Delta} \ar[u,"\Tr_f"']
\end{tikzcd}
\]
\end{booktheorem}

\begin{proof}[Proof (!)]
Show first that the analogous diagram of $f^{y}$, $\rho_f$ and V is valid for $f$ finite. Then follow the construction of $\Tr_f$ in section~\hyperref[sec:6.4]{4} to show that it is true in general.
\end{proof}\par\endgroup
